\scanpage{24}
\mspage{-2.23-}
% archive sheet number: 23

Therefore, in view of lemma 2.9, $\Delta(G)$ can be viewed as the result of pasting together spheres, one for each $T\subset G$, with sufficient identifications so that any two points are both on at least one sphere. The subsets $\Delta(T)$ will be called \uline{skeletons} of $\Delta(G)$. We can push this structure further in 2 ways: first the spheres $\Delta(T)$ can be broken up by considering the function $P(\delta)$:

\uline{Definition} \uline{2.7}: For all parabolic subgroups $P$, let $\Delta_{\handwritten{\unsure{P|p}}}(G)$ be the set of $\delta\in\Delta(G)$ such that $P(\delta)\supset P$.

To describe $\Delta_P(G)$ explicitly, let $T$ be a maximal torus of $P$. Then $\Delta_P(G)\subset\Delta(T)$ by lemma 2.9. Suppose $P$ is spanned by $T$ and by additive subgroups $\handwritten{\mathbb{G}}_a\subset P$ such that
\[
 \alpha^{-1}\cdot\gamma\cdot\alpha=\handwritten{\chi}(\alpha)\cdot\gamma
\]
for $\alpha\in T_k$, $\gamma\in(\handwritten{\mathbb{G}}_a)_k$, and a suitable character $\handwritten{\chi}$ of $T$. Then one checks that if $\delta\in\Delta(T)$ is induced by $\bar{\lambda}\in\unsure{\Gamma|\Pi}^{\handwritten{\mathbb{R}}}(T)$, and if $\handwritten{\chi}$ induces the linear functional $\bar{\handwritten{\chi}}$ on $\Gamma^{\handwritten{\mathbb{R}}}(T)$, then
\[
 P(\delta)\supset P\quad\handwritten{\Longleftrightarrow}\quad\handwritten{\forall}\text{ such }\handwritten{\chi},\quad\bar{\handwritten{\chi}}(\bar{\lambda})\leq 0.
\]
Therefore, $\Delta_P(G)$ is simple the set of rational points on an intersection of hemi-spheres, (each determined by a rational linear functional, in fact). If $G$ is actually semi-simple, it can be shown that $\Delta_{\handwritten{P}}(G)$ is even the set of rational points on a spherical simplex, and hence the collection of $\Delta_{\handwritten{P}}(G)$ constitutes a ``triangulation'', and $\Delta(G)$ is a complex formed from the simplices $\Delta_{\handwritten{P}}(G)$. This is the motivation
